% Template for post or presentation abstracts submitted to ILAS 2022
% 1. Fill in the presenter's information (name, affiliation, email
% address) and talk information (title of presentation, abstract).
% 2. Compile to make sure there are no errors.
% 3. Save the .tex file for your abstract with a distinctive name,
% ideally "FamilynameGivenname.tex".
% 4. Upload your .tex file to https://clr.ie/129557
%       
% * Please keep your LaTeX commands simple, and avoid the use of
%    user-defined macros.
% * Include details of co-authors and funding acknowledgements after the abstract.
% * Upload only the LaTeX file (with .tex extension).
% * Questions: Contact Niall Madden (Niall.Madden@NUIGalway.ie)

\documentclass[a4paper]{amsart} 
\usepackage{amssymb} 
\usepackage{hyperref}
\pagestyle{empty}
\date{} 

%% IGNORE THE NEXT 4 LINES
\makeatletter % 
\def\labelenumi{\theenumi.}
\def\theenumi{\@arabic\c@enumi}
\makeatother
%% STOP IGNORING NOW

\begin{document}

\title{Accurate Bidiagonal Decompositions of Structured Totally Nonnegative Matrices with Repeated Nodes}
\author{Plamen Koev} %Presenting author only
\address{San Jose State University} % Affiliation of presenting author
\email{plamen.koev@sjsu.edu} % Email address of presenting author only

\maketitle

% The abstract  ***no more than one page***

The decomposition of a totally nonnegative matrix (one all of whose minors are nonnegative) as a product of nonnegative bidiagonals is a powerful tool for studying the properties of these matrices \cite{fallat01} as well as performing numerical computations to high relative accuracy \cite{koevSTN}.

The conventional bidiagonal decompositions resulting from the complete Neville elimination, when applied to totally nonnegative matrices of the Vandermonde or Cauchy type have singularities that are when some of the nodes defining those matrices coincide.

For example, the bidiagonal decomposition of a $3\times 3$ Vandermonde matrix is:
$$
\left[\begin{matrix} 
1 & x & x^2 \\
1 & y & y^2 \\
1 & z & z^2 
\end{matrix}\right]
\!\!=\!\!
\left[\begin{matrix}
1 \\
& 1\\
& 1 & 1
\end{matrix}\right]
\!\!\!\!
\left[\begin{matrix}
1 \\
1 & 1\\
& \frac{z-y}{y-x} & 1
\end{matrix}\right]
\!\!\!\!
\left[\begin{matrix}
1 \\
& \!\!y-x\\
&  &\!\!\!\!\! (z-x)(z-y)
\end{matrix}\right]
\!\!\!\!
\left[\begin{matrix}
1 & x\\
& 1 & y\\
&  & 1
\end{matrix}\right]
\!\!\!\!
\left[\begin{matrix}
1   \\
& 1& x\\
& & 1 
\end{matrix}\right],
$$
which is not defined when $x=y$. 

By relaxing the requirement for the bidiagonal factors to have unit diagonals (which has no detrimental effects on our ability to study or compute with these matrices), the singularity at $x=y$ can be removed:
$$
\left[\begin{matrix} 
1 & x & x^2 \\
1 & y & y^2 \\
1 & z & z^2 
\end{matrix}\right]
=
\left[\begin{matrix}
1 \\
& 1\\
& 1 & z-y
\end{matrix}\right]
\!\!\!\!
\left[\begin{matrix}
1 \\
1 & y-x \\
& 1 & z-x
\end{matrix}\right]
\!\!\!\!
\left[\begin{matrix}
1 \\
& 1 \\
&  & 1
\end{matrix}\right]
\!\!\!\!
\left[\begin{matrix}
1 & x\\
& 1 & y\\
&  & 1
\end{matrix}\right]
\!\!\!\!
\left[\begin{matrix}
1   \\
& 1& x\\
& & 1 
\end{matrix}\right].
$$

In this talk we will present the technique to systematically remove the singularities in the bidiagonal decompositions of many classes of structured totally nonnegative matrices with repeated nodes, such as (rational, h-, q-Bernstein-) Vandermonde, Lupas, Cauchy, Cauchy Vandermonde matrices, among others. 

Practical examples of computing with these matrices will also be presented.

%% Coauthor details here (optional - delete if not applicable)
% and funding acknowledgment (optional - delete if not applicable)
\emph{This is joint work with Jorge Delgado, Ana Marco, Jos\'e-Javier Mart\'nezd, Juan Manuel Pe\~na, Per-Olof Persson, and Steven Spasov.
Supported by the SJSU Woodward Fund.}

\begin{thebibliography}{1}

\bibitem{fallat01}
S.~M. Fallat.
\newblock Bidiagonal factorizations of totally nonnegative matrices.
\newblock {\em Amer. Math. Monthly}, 108(8):697--712, 2001.

\bibitem{koevSTN}
P.~Koev.
\newblock Accurate eigenvalues and zero {J}ordan blocks of (singular) totally
  nonnegative matrices.
\newblock {\em Numer.\ {M}ath.}, 141:693--713, 2019.
\end{thebibliography}


\end{document}


% Please leave the next bit alone  
%%% Local Variables: 
%%% mode: latex
%%% TeX-master: "../ILAS-2022"
%%% End:
